\section{Research Questions and Methodology}
\label{sec:methodology}

\subsection{Research questions}

We organize the evaluation around three questions.

\begin{description}[leftmargin=3.2em,style=nextline]
  \item[RQ1: Layering cost.] What latency and utilization change is introduced
        by a pass-through gateway relative to direct \vllm{}?
  \item[RQ2: Outer batching.] At matched offered load, does fixed gateway
        batching improve throughput or latency after accounting for queue time
        and the backend's own continuous batching?
  \item[RQ3: Adaptation.] Does a frozen low-load bypass policy avoid harmful
        waiting while preserving useful backlog coalescing?
\end{description}

For RQ2 and RQ3, the confirmatory comparison is paired by repetition seed.
The null is zero mean paired difference. We report effect direction and 95\%
Student-$t$ intervals over three run-level differences; we do not label a
difference statistically significant from request counts within a run.

\subsection{Testbed}

All measurements run locally in WSL2 Ubuntu on a Windows host with an
\gpu{}. The backend is \vllm{} 0.29.0 in eager mode with float16 weights,
maximum model length 2,048, GPU-memory utilization 0.85, and at most eight
sequences. The model is \model{} at immutable revision
\texttt{989aa7980e4cf806f80c7fef2b1adb7bc71aa306}
\citep{qwen2025technical}. Generation is deterministic
($T=0$) with at most 64 output tokens. Direct traffic targets port 8001;
gateway modes target \adip{} on port 8000, which forwards to the same backend.

\begin{table}[t]
  \centering
  \small
  \caption{Frozen primary experiment. Every cell has three run-level
  repetitions.}
  \label{tab:experiment}
  \begin{tabular}{ll}
    \toprule
    Factor & Value \\
    \midrule
    GPU & NVIDIA RTX 3060, 12\,GB \\
    Model & Qwen2.5-1.5B-Instruct, float16 \\
    Backend & vLLM 0.29.0, eager execution \\
    Modes & direct, pass-through, fixed, adaptive \\
    Offered rates & 0.5, 1.0, 1.5, 1.8 requests/s \\
    Prompt bucket & 8--128 tokenizer-derived input tokens \\
    Output cap & 64 tokens \\
    Warmup & 50 sequential requests per condition \\
    Measurement & 120 seconds, open loop \\
    Repetitions & seeds 1729, 2718, 3141 \\
    Fixed policy & max batch 8, max wait 1\,ms \\
    Adaptive policy & max batch 8, max wait 20\,ms \\
    \bottomrule
  \end{tabular}
\end{table}

\subsection{Workload and controls}

The workload is synthetic and length controlled, not a production trace. Four
project-authored seed passages are expanded deterministically until they fall
inside the 8--128 token bucket under the exact model tokenizer. All requests
use the same model, output cap, and temperature, making them compatible for
gateway batching. The benchmark generates the complete constant-rate arrival
schedule before execution. Each request sleeps until an absolute target time;
later requests do not wait for earlier completions. This open-loop design
avoids coordinated omission when the system slows down.

Fifty warmup requests precede each measured interval. Warmups are excluded
from every aggregate. A five-second cooldown separates conditions. The final
matrix contains $4\times4\times3=48$ conditions and 5,760 seconds of measured
traffic, excluding warmup and cooldown. The fixed 1\,ms comparator and adaptive
parameters were frozen from calibration and pilot data before the final matrix.
The runner executes the declared order by mode (direct, pass-through, fixed,
adaptive), then repetition, then increasing rate. This deterministic order
makes interruption recovery auditable but does not balance slow temporal or
thermal drift. Seed pairing controls workload construction, not wall-clock
position; we treat the lack of randomized interleaving as a validity
limitation rather than implying otherwise.

\subsection{Validity gate and failure accounting}

Host scheduling can invalidate an open-loop experiment even when the server
returns success. We therefore record target and actual arrival time for every
request. A condition is marked \texttt{invalid\_harness} when more than 1\%
of measured requests arrive over 50\,ms late. The raw records remain stored,
but the condition is excluded from performance claims. Ordinary incomplete
directories are rejected by the resume logic rather than appended to. Request
failures remain in the denominator and are never converted to zero latency.

\subsection{Metrics and uncertainty}

The request is the latency observation unit; the run is the replication unit.
For each run we compute achieved requests/s, output tokens/s, success and
failure counts, p50/p95/p99 end-to-end latency, mean and p95 gateway queue time,
mean backend time, mean batch size, GPU utilization, peak VRAM, power, and p95
arrival drift. We then average the three run-level values for each condition
and report a two-sided 95\% Student-$t$ interval. With only three repetitions,
these intervals are necessarily wide; figures show each run point to prevent
the mean from hiding instability.

Each request record carries its outer batch ID and reported batch size. A
request-wise average of that field is size biased: a group of four contributes
four observations while a singleton contributes one. For mechanism claims we
therefore group records by batch ID and report the call-weighted mean number of
prompts per backend HTTP call. The raw request-weighted mean remains in each
run summary for auditability but is not labeled as a call-level average. The
same unique-call grouping yields $\bar S_{\mathrm{call}}$ for
Eq.~\ref{eq:occupancy-prediction}; the analyzer reports measured group size,
predicted group size, and their ratio for every fixed and adaptive run.

The primary effects are paired by rate and seed:
\begin{align}
  \Delta_{\mathrm{proxy}} &= M_{\mathrm{pass}}-M_{\mathrm{direct}},\\
  \Delta_{\mathrm{fixed}} &= M_{\mathrm{fixed}}-M_{\mathrm{pass}},\\
  \Delta_{\mathrm{adaptive}} &= M_{\mathrm{adaptive}}-M_{\mathrm{fixed}},
\end{align}
for metric $M$. Absolute and relative deltas are both retained. We interpret
throughput separately from goodput and offered load separately from achieved
load.
